% GENERATED FILE. Edit canonical lessons, then run npm run generate:books.
% canonical-source-hash: fnv1a64:370e484c725d40aa
% canonical-lessons: FR-C212-fixer, FR-C212-freiner, FR-C212-fonctionner, FR-C212-flotter, FR-C212-poursuivre

\chapter{To fix, To brake, To function, To float, To chase}
\label{ch:fr-to-fix-to-brake-to-function-to-float-to-chase}
\hlchaptermodality{212}

\begin{chapteropening}
\textbf{By the end of this chapter:} \emph{I can say to fix, to brake, to function, to float and to chase.}

Complete the last lesson of chapter 212: I can say to fix, to brake, to function, to float and to chase.
\end{chapteropening}

\section[fixer]{fixer — to fix, to stare at}

\label{lesson:FR-C212-fixer}



\begin{quote}
Before the new one: say the French for to transport, then the French for to kill.
\end{quote}

\subsection*{You'll want to know: fixer}
\textbf{fixer} — \textquotedblleft{}to fix, to stare at\textquotedblright{}.

\begin{cousinweb}
\textbf{fixer} — built from the word for \textquotedblleft{}fixed\textquotedblright{}, $\leftarrow$ Latin \emph{fixus}, from \emph{figere}, \textquotedblleft{}to fasten\textquotedblright{}. (EtymDB 2.1)
\end{cousinweb}

\begin{grammarlens}[title={What you've built}]
One more word for this chapter.
\end{grammarlens}

\subsection*{Your turn}
\begin{itemize}
\raggedright
  \item \emph{Say it:} \emph{fixer}
  \item \emph{Say it:} \emph{fixer}, once more
  \item \emph{Say it:} say \emph{fixer}
  \item \emph{Recall:} say the French for to transport, then the French for to kill, then say \emph{fixer} again
\end{itemize}

\begin{tcolorbox}[breakable,colback=teal!4,colframe=teal!35!black,title={Before you move on}]
What does fixer mean? (\textquotedblleft{}To fix, to stare at\textquotedblright{}.) Say it once more.
\end{tcolorbox}
\section[freiner]{freiner — to brake}

\label{lesson:FR-C212-freiner}



\begin{quote}
Before the new one: say the French for to kill, then the French for to fix.
\end{quote}

\subsection*{You'll want to know: freiner}
\textbf{freiner} — \textquotedblleft{}to brake\textquotedblright{}.

\begin{cousinweb}
\textbf{freiner} — built from the word for \textquotedblleft{}a bit, a brake\textquotedblright{}, $\leftarrow$ Latin \emph{frenum}, \textquotedblleft{}a bridle\textquotedblright{} (Lewis \& Short).
\end{cousinweb}

\begin{grammarlens}[title={What you've built}]
One more word for this chapter.
\end{grammarlens}

\subsection*{Your turn}
\begin{itemize}
\raggedright
  \item \emph{Say it:} \emph{freiner}
  \item \emph{Say it:} \emph{freiner}, once more
  \item \emph{Say it:} say \emph{freiner}
  \item \emph{Recall:} say the French for to kill, then the French for to fix, then say \emph{freiner} again
\end{itemize}

\begin{tcolorbox}[breakable,colback=teal!4,colframe=teal!35!black,title={Before you move on}]
What does freiner mean? (\textquotedblleft{}To brake\textquotedblright{}.) Say it once more.
\end{tcolorbox}
\section[fonctionner]{fonctionner — to function, to work}

\label{lesson:FR-C212-fonctionner}



\begin{quote}
Before the new one: say the French for to fix, then the French for to brake.
\end{quote}

\subsection*{You'll want to know: fonctionner}
\textbf{fonctionner} — \textquotedblleft{}to function, to work\textquotedblright{}.

\begin{cousinweb}
\textbf{fonctionner} — built from the word for \textquotedblleft{}function\textquotedblright{}, $\leftarrow$ Latin \emph{functio}, \textquotedblleft{}a performing\textquotedblright{}, with \emph{-er}. (EtymDB 2.1)
\end{cousinweb}

\begin{grammarlens}[title={What you've built}]
One more word for this chapter.
\end{grammarlens}

\subsection*{Your turn}
\begin{itemize}
\raggedright
  \item \emph{Say it:} \emph{fonctionner}
  \item \emph{Say it:} \emph{fonctionner}, once more
  \item \emph{Say it:} say \emph{fonctionner}
  \item \emph{Recall:} say the French for to fix, then the French for to brake, then say \emph{fonctionner} again
\end{itemize}

\begin{tcolorbox}[breakable,colback=teal!4,colframe=teal!35!black,title={Before you move on}]
What does fonctionner mean? (\textquotedblleft{}To function, to work\textquotedblright{}.) Say it once more.
\end{tcolorbox}
\section[flotter]{flotter — to float}

\label{lesson:FR-C212-flotter}



\begin{quote}
Before the new one: say the French for to brake, then the French for to function.
\end{quote}

\subsection*{You'll want to know: flotter}
\textbf{flotter} — \textquotedblleft{}to float\textquotedblright{}.

\begin{cousinweb}
\textbf{flotter} $\leftarrow$ Old French \emph{floter} $\leftarrow$ Frankish \emph{*flotan}, \textquotedblleft{}to float\textquotedblright{} — borrowed. (EtymDB 2.1)
\end{cousinweb}

\begin{grammarlens}[title={What you've built}]
One more word for this chapter.
\end{grammarlens}

\subsection*{Your turn}
\begin{itemize}
\raggedright
  \item \emph{Say it:} \emph{flotter}
  \item \emph{Say it:} \emph{flotter}, once more
  \item \emph{Say it:} say \emph{flotter}
  \item \emph{Recall:} say the French for to brake, then the French for to function, then say \emph{flotter} again
\end{itemize}

\begin{tcolorbox}[breakable,colback=teal!4,colframe=teal!35!black,title={Before you move on}]
What does flotter mean? (\textquotedblleft{}To float\textquotedblright{}.) Say it once more.
\end{tcolorbox}
\section[poursuivre]{poursuivre — to chase}

\label{lesson:FR-C212-poursuivre}



\begin{quote}
Before the new one: say the French for to function, then the French for to float.
\end{quote}

\subsection*{You'll want to know: poursuivre}
\textbf{poursuivre} — \textquotedblleft{}to chase\textquotedblright{}.

\begin{cousinweb}
\textbf{poursuivre} $\leftarrow$ Old French \emph{porsivre} (EtymDB 2.1), from \emph{pour-}, \textquotedblleft{}forward\textquotedblright{}, and the verb for \textquotedblleft{}to follow\textquotedblright{}, $\leftarrow$ Latin \emph{sequi} — inherited.
\end{cousinweb}

\begin{grammarlens}[title={What you've built}]
One more word for this chapter.
\end{grammarlens}

\subsection*{Your turn}
\begin{itemize}
\raggedright
  \item \emph{Say it:} \emph{poursuivre}
  \item \emph{Say it:} \emph{poursuivre}, once more
  \item \emph{Say it:} say \emph{poursuivre}
  \item \emph{Recall:} say the French for to function, then the French for to float, then say \emph{poursuivre} again
\end{itemize}

\begin{tcolorbox}[breakable,colback=teal!4,colframe=teal!35!black,title={Before you move on}]
What does poursuivre mean? (\textquotedblleft{}To chase\textquotedblright{}.) Say it once more.
\end{tcolorbox}
